On the CR5, I aligned observations to RGB timestamps and carried policy chunks through WebSocket serving to robot execution. In the HERO RoboMaster team, I worked on ROS2 messaging, host-MCU communication, and time synchronization. These experiences taught me to diagnose individual handoffs, but I still need to understand how sensor input, cloud deployment, and computing hardware jointly shape learning-system performance. CityUHK's MSc in Computer and Information Engineering maps directly onto that cross-layer gap.

With the annual availability of all three technical electives still to be confirmed, EE5811 Topics in Computer Vision and EE5438 Applied Deep Learning would turn sensor data into learned interpretation. EE5608 Modern Cloud Architecture and Deployment Practices, if available and selected as one core choice, would carry model outputs across service boundaries; EE6625 Hardware Architectures for Artificial Intelligence and Machine Learning would expose the computing substrate beneath inference. I would follow one synchronized observation through a cloud-deployed vision service, measure how deployment and hardware choices affect response time, and connect that response to the robot command. This cross-layer view would help me design embodied systems in which computing architecture supports the behavior required at the machine.
